\centerline{\bf \Huge Текстовые задачи на движение F}

\centerline{\bf \Huge Движение по прямой дороге}

\problem{\textit{\textbf{1}}} Из городов А и В, расстояние между которыми равно 330 км, навстречу друг другу одновременно выехали два автомобиля и встретились через 3 часа на расстоянии 180 км от города В. Найдите скорость автомобиля, выехавшего из города А. Ответ дайте в км/ч.

\problem{\textit{\textbf{2}}} Два пешехода АА и РВ отправляются одновременно в одном направлении из одного и того же места на прогулку по аллее парка. Скорость АА на 0,5 км/ч больше скорости РВ. Через сколько минут расстояние между пешеходами станет равным 425 метрам?

\problem{\textit{\textbf{3}}} Из пункта А в пункт В, расстояние между которыми 50 км, одновременно выехали автомобилист и велосипедист. Известно, что в час автомобилист проезжает на 65 км больше, чем велосипедист. Определите скорость велосипедиста, если известно, что он прибыл в пункт В на 4 часа 20 минут позже автомобилиста. Ответ дайте в км/ч.

\problem{\textit{\textbf{4}}} Велосипедист выехал с постоянной скоростью из города А в город В, расстояние между которыми равно 180 км. На следующий день он отправился обратно в А со скоростью на 8 км/ч больше прежней. По дороге он сделал остановку на 8 часов. В результате велосипедист затратил на обратный путь столько же времени, сколько на путь из А в В. Найдите скорость велосипедиста на пути из В в А. Ответ дайте в км/ч.

\problem{\textit{\textbf{5}}} Поезд, пройдя 450 км, был остановлен из-за снежного заноса. Через полчаса путь был расчищен, и машинист, увеличив скорость поезда на 15 км/ч, привел его на станцию без опоздания. Найдите первоначальную скорость поезда, если путь, пройденный им до остановки, составил 75\% всего пути.

\problem{\textit{\textbf{6}}} Два велосипедиста одновременно отправились в 224 километровый пробег. Первый ехал со скоростью, на 2 км/ч большей, чем скорость второго, и прибыл к финишу на 2 часа раньше второго. Найти скорость велосипедиста, пришедшего к финишу вторым. Ответ дайте в км/ч.

\problem{\textit{\textbf{7}}} Из городов А и В одновременно навстречу друг другу выехали мотоциклист и велосипедист. Мотоциклист приехал в В на 10 часов раньше, чем велосипедист приехал в А, а встретились они через 55 минут после выезда. Сколько часов затратил на путь из В в А велосипедист?

\newpage
\hspace{3mm}

\problem{\textit{\textbf{8}}} Товарный поезд каждую минуту проезжает на 300 метров меньше, чем скорый, и на путь в 420 км тратит времени на 3 часа больше, чем скорый. Найдите скорость товарного поезда. Ответ дайте в км/ч.

\problem{\textit{\textbf{9}}} Расстояние между городами А и В равно 348 км. Из города А в город В выехал автомобиль, а через 1 час следом за ним со скоростью 85 км/ч выехал мотоциклист, догнал автомобиль в городе С и повернул обратно. Когда он вернулся в А, автомобиль прибыл в В. Найдите расстояние от A до C . Ответ дайте в километрах.

\problem{\textit{\textbf{10}}} Из А в В одновременно выехали два автомобилиста. Первый проехал с постоянной скоростью весь путь. Второй проехал первую половину пути со скоростью, меньшей скорости первого на 12 км/ч, а вторую половину пути - со скоростью 90 км/ч, в результате чего прибыл в В одновременно с первым автомобилистом. Найдите скорость первого автомобилиста, если известно, что она больше 50 км/ч. Ответ дайте в км/ч.


\centerline{\bf \Huge Движение по замкнутой дороге}


\problem{\textit{\textbf{11}}} Два мотоциклиста стартуют одновременно в одном направлении из двух диаметрально противоположных точек круговой трассы, длина которой равна 22 км. Через сколько минут мотоциклисты поравняются в первый раз, если скорость одного из них на 20 км/ч больше скорости другого?

\problem{\textit{\textbf{12}}} Из одной точки круговой трассы, длина которой равна 12 км, одновременно в одном направлении стартовали два автомобиля. Скорость первого автомобиля равна 101 км/ч, и через 20 минут после старта он опережал второй автомобиль на один круг. Найдите скорость второго автомобиля. Ответ дайте в км/ч.

\problem{\textit{\textbf{13}}} Из пункта А круговой трассы выехал велосипедист, а через 40 минут следом за ним отправился мотоциклист. Через 8 минут после отправления он догнал велосипедиста в первый раз, а еще через 36 минут после этого догнал его во второй раз. Найдите скорость мотоциклиста, если длина трассы равна 30 км. Ответ дайте в км/ч.

\problem{\textit{\textbf{14}}} Часы со стрелками показывают 4 часа 45 минут. Через сколько минут минутная стрелка в седьмой раз поравняется с часовой?

\problem{\textit{\textbf{15}}} На соревнованиях по кольцевой трассе один лыжник проходит круг на 2 мин быстрее другого и через час обошел его ровно на круг. За какое время каждый лыжник проходит круг?

\newpage

\centerline{\bf \Huge Движение по реке}


\problem{\textit{\textbf{16}}} Моторная лодка прошла против течения реки 195 км и вернулась в пункт отправления, затратив на обратный путь на 2 часа меньше. Найдите скорость течения, если скорость лодки в неподвижной воде равна 14 км/ч. Ответ дайте в км/ч.

\problem{\textit{\textbf{17}}} Моторная лодка прошла против течения реки 221 км и вернулась в пункт отправления, затратив на обратный путь на 4 часа меньше. Найдите скорость лодки в неподвижной воде, если скорость течения равна 2 км/ч. Ответ дайте в км/ч.

\problem{\textit{\textbf{18}}} Катер в 11:00 вышел из пункта А в пункт В, расположенный в 30 км от А. Пробыв в пункте В 2 часа 40 минут, катер отправился назад и вернулся в пункт А в 19:00 того же дня. Определите (в км/ч) собственную скорость катера, если известно, что скорость течения реки 3 км/ч.

\problem{\textit{\textbf{19}}} От пристани А к пристани В, расстояние между которыми равно 168 км, отправился с постоянной скоростью первый теплоход, а через 2 часа после этого следом за ним, со скоростью на 2 км/ч большей, отправился второй. Найдите скорость первого теплохода, если в пункт В оба теплохода прибыли одновременно. Ответ дайте в км/ч.

\problem{\textit{\textbf{20}}} Расстояние между пристанями А и В равно 140 км. Из А в В по течению реки отправился плот, а через 3 часа вслед за ним отправилась яхта, которая, прибыв в пункт В, тотчас повернула обратно и возвратилась в А. К этому времени плот прошел 60 км. Найдите скорость яхты в неподвижной воде, если скорость течения реки равна 3 км/ч. Ответ дайте в км/ч.

\problem{\textit{\textbf{21}}} Пристани A и B расположены на озере, расстояние между ними равно 416 км. Баржа отправилась с постоянной скоростью из А в В. На следующий день после прибытия она отправилась обратно со скоростью на 3 км/ч больше прежней, сделав по пути остановку на 6 часов. В результате она затратила на обратный путь столько же времени, сколько на путь из А в В. Найдите скорость баржи на пути из А в В. Ответ дайте в км/ч.


\centerline{\bf \Huge Движение протяженных тел}


\problem{\textit{\textbf{22}}} Поезд, двигаясь равномерно со скоростью 50 км/ч, проезжает мимо придорожного столба за 72 секунды. Найдите длину поезда в метрах.

\problem{\textit{\textbf{23}}} Поезд, двигаясь равномерно со скоростью 70 км/ч, проезжает мимо лесополосы, длина которой равна 1000 метров, за 1 минуту 48 секунд. Найдите длину поезда в метрах.

\newpage
\hspace{3mm}

\problem{\textit{\textbf{24}}} По двум параллельным железнодорожным путям в одном направлении следуют пассажирский и товарный поезда, скорости которых равны соответственно 80 км/ч и 50 км/ч. Длина товарного поезда равна 800 метрам. Найдите длину пассажирского поезда, если время, за которое он прошел мимо товарного поезда, равно 2 минутам. Ответ дайте в метрах.

\problem{\textit{\textbf{25}}} По морю параллельными курсами в одном направлении следуют два сухогруза: первый длиной 130 метров, второй - длиной 120 метров. Сначала второй сухогруз отстает от первого, и в некоторый момент времени расстояние от кормы первого сухогруза до носа второго составляет 600 метров. Через 11 минут после этого уже первый сухогруз отстает от второго так, что расстояние от кормы второго сухогруза до носа первого равно 800 метрам. На сколько километров в час скорость первого сухогруза меньше скорости второго?


\centerline{\bf \Huge Средняя скорость}


\problem{\textit{\textbf{26}}} Первые три часа автомобиль ехал со скоростью 70 км/ч, следующий час - со скоростью 65 км/ч, а затем один час - со скоростью 45 км/ч. Найдите среднюю скорость автомобиля на протяжении всего пути. Ответ дайте в км/ч.

\problem{\textit{\textbf{27}}} Первые 120 км автомобиль ехал со скоростью 50 км/ч, следующие 160 км - со скоростью 100 км/ч, а затем 120 км - со скоростью 120 км/ч. Найдите среднюю скорость автомобиля на протяжении всего пути. Ответ дайте в км/ч.

\problem{\textit{\textbf{28}}} Путешественник переплыл море на яхте со средней скоростью 28 км/ч. Обратно он летел на спортивном самолете со скоростью 532 км/ч. Найдите среднюю скорость путешественника на протяжении всего пути. Ответ дайте в км/ч.

\problem{\textit{\textbf{29}}} Первую треть трассы автомобиль ехал со скоростью 100 км/ч, вторую треть - со скоростью 75 км/ч, а последнюю - со скоростью 60 км/ч. Найдите среднюю скорость автомобиля на протяжении всего пути. Ответ дайте в км/ч.

\problem{\textit{\textbf{30}}} Половину времени, затраченного на дорогу, автомобиль ехал со скоростью 67 км/ч, а вторую половину времени - со скоростью 79 км/ч. Найдите среднюю скорость автомобиля на протяжении всего пути. Ответ дайте в км/ч.
